\documentclass[aspectratio=169]{beamer}

% Use the UFS theme
\usetheme{UFS}

% Additional packages
\usepackage[utf8]{inputenc}
\usepackage[portuguese]{babel}
\usepackage{amsmath,amssymb}
\usepackage{graphicx}
\usepackage{booktabs}
\usepackage{listings}
\usepackage{xcolor}
\usepackage{tikz}

% --- Paleta VS Code Light+ ---
\definecolor{bgcolor}{HTML}{FFFFFF}       % fundo branco
\definecolor{linecolor}{HTML}{DDDDDD}     % borda cinza clara
\definecolor{numcolor}{HTML}{999999}      % números de linha
\definecolor{kwcolor}{HTML}{0000FF}       % keywords: azul
\definecolor{strcolor}{HTML}{A31515}      % strings: vermelho escuro
\definecolor{cmtcolor}{HTML}{008000}      % comentários: verde
\definecolor{fncolor}{HTML}{795E26}       % funções: dourado/marrom
\definecolor{typecolor}{HTML}{267F99}     % tipos/classes: azul-petróleo

\lstdefinestyle{modernstyle}{
    backgroundcolor=\color{bgcolor},
    basicstyle=\ttfamily\footnotesize\color{black},
    keywordstyle=\color{kwcolor}\bfseries,
    stringstyle=\color{strcolor},
    commentstyle=\color{cmtcolor}\itshape,
    numberstyle=\tiny\color{numcolor},
    % --- Layout ---
    numbers=left,
    numbersep=12pt,
    xleftmargin=20pt,
    framexleftmargin=20pt,
    % --- Quebra de linha ---
    breaklines=true,
    breakatwhitespace=false,
    % --- Espaçamento e tabs ---
    tabsize=4,
    keepspaces=true,
    showspaces=false,
    showstringspaces=false,
    showtabs=false,
    % --- Moldura ---
    frame=single,
    rulecolor=\color{linecolor},
    framesep=4pt,
    % --- Legenda ---
    captionpos=b,
    upquote=true,
    columns=fullflexible,
}
\lstset{style=modernstyle}

% --- Metadados ---
\title{COMP0496 -- Programação A \\ Programação Orientada a Objetos --- Parte 1}
\subtitle{Classes, Objetos, Atributos e Métodos}
\author{Prof.\ Dr.\ André Yoshiaki Kashiwabara}
\institute{DCOMP -- Universidade Federal de Sergipe\\
\texttt{andreyoshiaki@dcomp.ufs.br}}
\date{}


\begin{document}

% ============================================================
% SLIDE 1 – Capa
% ============================================================
\begin{frame}[plain]
  \titlepage
\end{frame}

% ============================================================
% SLIDE 2 – Roteiro
% ============================================================
\begin{frame}{Roteiro}
  \tableofcontents
\end{frame}

% ############################################################
% SEÇÃO 1: Aquecimento — Tudo em Python é Objeto
% ############################################################
\section{Aquecimento --- Tudo em Python é Objeto}

% ============================================================
% SLIDE 3
% ============================================================
\begin{frame}[fragile]{Tudo em Python é Objeto}
  Você já usa objetos sem saber! Veja o que \texttt{type()} retorna:

  \begin{lstlisting}[language=Python]
>>> type("Oi")
<class 'str'>
>>> type([1, 2, 3])
<class 'list'>
>>> type(42)
<class 'int'>
  \end{lstlisting}

  \vspace{0.3cm}
  Cada valor em Python pertence a uma \textbf{classe} --- um tipo que define
  quais operações são possíveis.
\end{frame}

% ============================================================
% SLIDE 4
% ============================================================
\begin{frame}[fragile]{Objetos Possuem Métodos}
  Métodos são funções que ``pertencem'' a um objeto:

  \begin{lstlisting}[language=Python]
>>> "programacao".upper()
'PROGRAMACAO'

>>> frutas = ["banana", "manga"]
>>> frutas.append("uva")
>>> frutas
['banana', 'manga', 'uva']
  \end{lstlisting}

  \vspace{0.3cm}
  \begin{block}{Ponto-chave}
    POO = criar \textbf{seus próprios tipos} (classes)
    com seus próprios métodos e atributos.
  \end{block}
\end{frame}

% ############################################################
% SEÇÃO 2: O Problema — A Sopa de Variáveis
% ############################################################
\section{O Problema --- A Sopa de Variáveis}

% ============================================================
% SLIDE 5
% ============================================================
\begin{frame}[fragile]{O Cenário: Lanchonete ``O Podrão''}
  Você abriu um food truck e precisa de um sistema para gerenciar
  o cardápio.

  \vspace{0.3cm}
  Primeira tentativa: variáveis soltas.

  \begin{lstlisting}[language=Python]
lanche1_nome = "X-Tudo"
lanche1_preco = 18.50
lanche1_ingredientes = "pao, hamburguer, queijo, ovo, bacon"

lanche2_nome = "X-Salada"
lanche2_preco = 14.00
lanche2_ingredientes = "pao, hamburguer, queijo, alface, tomate"
  \end{lstlisting}
\end{frame}

% ============================================================
% SLIDE 6
% ============================================================
\begin{frame}{Por que Isso é Ruim?}
  \begin{itemize}
    \item \textbf{Dados soltos}: nome, preço e ingredientes não estão
          conectados --- qualquer variável pode ser alterada isoladamente
    \item \textbf{Não escala}: com 50 lanches, seriam 150 variáveis!
    \item \textbf{Erros de digitação}: \texttt{lanche1\_preco} vs.\
          \texttt{lanche1\_prço} --- sem aviso
    \item \textbf{Sem comportamento}: não há como pedir a um lanche
          que ``se descreva''
  \end{itemize}

  \vspace{0.5cm}
  \begin{alertblock}{Analogia}
    É como anotar pedidos em guardanapos separados:
    fácil perder, fácil trocar, impossível de organizar.
  \end{alertblock}
\end{frame}

% ############################################################
% SEÇÃO 3: A Solução — Classes e Objetos
% ############################################################
\section{A Solução --- Classes e Objetos}

% ============================================================
% SLIDE 7
% ============================================================
\begin{frame}{Receita vs.\ Lanche Pronto}
  \begin{columns}[T]
    \begin{column}{0.48\textwidth}
      \textbf{Classe = Receita}
      \begin{itemize}
        \item Define \emph{quais} informações um lanche tem
        \item Define \emph{o que} um lanche pode fazer
        \item É um molde, um modelo
        \item Existe uma única vez no código
      \end{itemize}
    \end{column}
    \begin{column}{0.48\textwidth}
      \textbf{Objeto = Lanche Pronto}
      \begin{itemize}
        \item Preenche a receita com dados reais
        \item Cada objeto é independente
        \item Pode haver muitos objetos da mesma classe
        \item Vive na memória do programa
      \end{itemize}
    \end{column}
  \end{columns}
\end{frame}

% ============================================================
% SLIDE 8
% ============================================================
\begin{frame}[fragile]{Primeira Classe: \texttt{Lanche}}
  \begin{lstlisting}[language=Python]
class Lanche:
    """Representa um lanche do cardapio."""

    def __init__(self, nome, preco, ingredientes):
        self.nome = nome
        self.preco = preco
        self.ingredientes = ingredientes

    def descrever(self):
        return f"{self.nome} - R${self.preco:.2f}"
  \end{lstlisting}
\end{frame}

% ============================================================
% SLIDE 9
% ============================================================
\begin{frame}{Anatomia do \texttt{\_\_init\_\_}}
  \begin{itemize}
    \item \texttt{\_\_init\_\_} é o \textbf{construtor} --- é chamado
          automaticamente ao criar um objeto
    \item \texttt{self} é uma referência ao \textbf{próprio objeto}
          que está sendo criado
    \item \texttt{self.nome = nome} cria um \textbf{atributo de instância}
          --- cada objeto terá o seu
  \end{itemize}

  \vspace{0.3cm}
  \begin{block}{Atenção}
    \texttt{self} \textbf{sempre} é o primeiro parâmetro de qualquer método,
    mas \textbf{não} é passado na chamada --- Python cuida disso.
  \end{block}
\end{frame}

% ============================================================
% SLIDE 10
% ============================================================
\begin{frame}[fragile]{Criando Objetos (Instâncias)}
  \begin{lstlisting}[language=Python]
xtudo = Lanche("X-Tudo", 18.50, "pao, hamburguer, queijo, ovo, bacon")
xsalada = Lanche("X-Salada", 14.00, "pao, hamburguer, queijo, alface")

print(xtudo.descrever())    # X-Tudo - R$18.50
print(xsalada.nome)         # X-Salada
print(xsalada.preco)        # 14.0
  \end{lstlisting}

  \vspace{0.3cm}
  Cada objeto é independente: alterar \texttt{xtudo.preco} não afeta
  \texttt{xsalada.preco}.
\end{frame}

% ============================================================
% SLIDE 11
% ============================================================
\begin{frame}[fragile]{Comparação: Antes vs.\ Depois}
  \begin{columns}[T]
    \begin{column}{0.48\textwidth}
      \textbf{Sem OOP}
      \begin{lstlisting}[language=Python, basicstyle=\ttfamily\scriptsize]
lanche1_nome = "X-Tudo"
lanche1_preco = 18.50
# ... e mais 148 variaveis
      \end{lstlisting}
    \end{column}
    \begin{column}{0.48\textwidth}
      \textbf{Com OOP}
      \begin{lstlisting}[language=Python, basicstyle=\ttfamily\scriptsize]
xtudo = Lanche("X-Tudo", 18.50,
               "pao, hamburguer...")
# 1 variavel = 1 lanche completo
      \end{lstlisting}
    \end{column}
  \end{columns}

  \vspace{0.5cm}
  Dados e comportamento ficam \textbf{juntos} --- é isso que
  ``orientação a objetos'' significa.
\end{frame}

% ============================================================
% SLIDE 12
% ============================================================
\begin{frame}{Vocabulário Essencial}
  \begin{table}
    \centering
    \begin{tabular}{ll}
      \toprule
      \textbf{Conceito} & \textbf{Significado} \\
      \midrule
      Classe         & Modelo/receita que define atributos e métodos \\
      Objeto         & Instância concreta de uma classe \\
      \texttt{\_\_init\_\_}  & Método construtor --- inicializa o objeto \\
      \texttt{self}  & Referência ao próprio objeto \\
      Atributo       & Variável que pertence ao objeto \\
      Método         & Função que pertence à classe \\
      \bottomrule
    \end{tabular}
  \end{table}
\end{frame}

% ============================================================
% SLIDE 13
% ============================================================
\begin{frame}[fragile]{Métodos: Adicionando Comportamento}
  \begin{lstlisting}[language=Python]
class Lanche:
    def __init__(self, nome, preco, ingredientes):
        self.nome = nome
        self.preco = preco
        self.ingredientes = ingredientes

    def descrever(self):
        return f"{self.nome} - R${self.preco:.2f}"

    def aplicar_desconto(self, percentual):
        self.preco *= (1 - percentual / 100)

    def tem_ingrediente(self, ingrediente):
        return ingrediente in self.ingredientes
  \end{lstlisting}
\end{frame}

% ============================================================
% SLIDE 14
% ============================================================
\begin{frame}[fragile]{Usando os Métodos}
  \begin{lstlisting}[language=Python]
xtudo = Lanche("X-Tudo", 18.50, "pao, hamburguer, queijo, ovo")

print(xtudo.descrever())             # X-Tudo - R$18.50

xtudo.aplicar_desconto(10)
print(xtudo.descrever())             # X-Tudo - R$16.65

print(xtudo.tem_ingrediente("ovo"))  # True
print(xtudo.tem_ingrediente("bacon"))# False
  \end{lstlisting}

  \vspace{0.3cm}
  O objeto \textbf{sabe} como se descrever, como aplicar desconto
  e como verificar ingredientes --- o comportamento está \emph{dentro} do objeto.
\end{frame}

% ############################################################
% SEÇÃO 4: Atributos de Classe vs Instância
% ############################################################
\section{Atributos de Classe vs Instância}

% ============================================================
% SLIDE 15
% ============================================================
\begin{frame}{Dois Tipos de Atributos}
  \begin{table}
    \centering
    \begin{tabular}{lll}
      \toprule
      & \textbf{Atributo de Classe} & \textbf{Atributo de Instância} \\
      \midrule
      Onde define  & Dentro da classe, fora de métodos & Dentro de \texttt{\_\_init\_\_} \\
      Compartilhado & Sim, por todos os objetos & Não, cada objeto tem o seu \\
      Acesso       & \texttt{Classe.atributo} & \texttt{self.atributo} \\
      Exemplo      & Taxa de serviço (10\%) & Nome do lanche \\
      \bottomrule
    \end{tabular}
  \end{table}
\end{frame}

% ============================================================
% SLIDE 16
% ============================================================
\begin{frame}[fragile]{Exemplo: Atributo de Classe}
  \begin{lstlisting}[language=Python]
class Lanche:
    taxa_servico = 0.10  # compartilhado por TODOS os lanches

    def __init__(self, nome, preco, ingredientes):
        self.nome = nome
        self.preco = preco
        self.ingredientes = ingredientes

    def preco_com_taxa(self):
        return self.preco * (1 + Lanche.taxa_servico)
  \end{lstlisting}

  \begin{lstlisting}[language=Python]
xtudo = Lanche("X-Tudo", 18.50, "pao, hamburguer, queijo")
print(xtudo.preco_com_taxa())  # 20.35

Lanche.taxa_servico = 0.15    # altera para TODOS
print(xtudo.preco_com_taxa())  # 21.275
  \end{lstlisting}
\end{frame}

% ============================================================
% SLIDE 17
% ============================================================
\begin{frame}[fragile]{Armadilha: Atributo de Classe Mutável}
  \begin{lstlisting}[language=Python]
class Lanche:
    extras = []  # PERIGO! Lista compartilhada

    def __init__(self, nome):
        self.nome = nome

    def adicionar_extra(self, extra):
        self.extras.append(extra)
  \end{lstlisting}

  \begin{lstlisting}[language=Python]
a = Lanche("X-Tudo")
b = Lanche("X-Salada")
a.adicionar_extra("bacon")
print(b.extras)  # ['bacon'] -- b tambem tem bacon?!
  \end{lstlisting}
\end{frame}

% ============================================================
% SLIDE 18
% ============================================================
\begin{frame}[fragile]{Correção: Criar Mutáveis no \texttt{\_\_init\_\_}}
  \begin{lstlisting}[language=Python]
class Lanche:
    def __init__(self, nome):
        self.nome = nome
        self.extras = []  # CORRETO: cada objeto tem sua lista

    def adicionar_extra(self, extra):
        self.extras.append(extra)
  \end{lstlisting}

  \begin{lstlisting}[language=Python]
a = Lanche("X-Tudo")
b = Lanche("X-Salada")
a.adicionar_extra("bacon")
print(a.extras)  # ['bacon']
print(b.extras)  # []  -- independente!
  \end{lstlisting}

  \begin{alertblock}{Regra de Ouro}
    Atributos mutáveis (listas, dicionários, conjuntos) devem ser
    criados em \texttt{\_\_init\_\_}, \textbf{nunca} como atributos de classe.
  \end{alertblock}
\end{frame}

% ############################################################
% SEÇÃO 5: Prática — Classe Cardapio
% ############################################################
\section{Prática --- Classe Cardápio}

% ============================================================
% SLIDE 19
% ============================================================
\begin{frame}{Composição: Objetos Dentro de Objetos}
  Um \texttt{Cardapio} contém uma \textbf{lista de objetos}
  \texttt{Lanche}.

  \vspace{0.3cm}
  Isso é chamado de \textbf{composição}: um objeto é composto
  por outros objetos.

  \vspace{0.3cm}
  \begin{block}{Analogia}
    O cardápio do food truck não é um lanche --- ele \emph{contém}
    lanches. A relação é ``tem-um'' (has-a), não ``é-um'' (is-a).
  \end{block}
\end{frame}

% ============================================================
% SLIDE 20
% ============================================================
\begin{frame}[fragile]{Classe \texttt{Cardapio}: Estrutura}
  \begin{lstlisting}[language=Python]
class Cardapio:
    """Gerencia a colecao de lanches do food truck."""

    def __init__(self, nome_food_truck):
        self.nome_food_truck = nome_food_truck
        self.lanches = []  # lista de objetos Lanche

    def adicionar(self, lanche):
        self.lanches.append(lanche)

    def mostrar(self):
        print(f"=== Cardapio: {self.nome_food_truck} ===")
        for lanche in self.lanches:
            print(f"  {lanche.descrever()}")
  \end{lstlisting}
\end{frame}

% ============================================================
% SLIDE 21
% ============================================================
\begin{frame}[fragile]{Classe \texttt{Cardapio}: Métodos de Consulta}
  \begin{lstlisting}[language=Python]
class Cardapio:
    # ... (continuacao)

    def preco_medio(self):
        if not self.lanches:
            return 0
        total = sum(l.preco for l in self.lanches)
        return total / len(self.lanches)

    def mais_caro(self):
        if not self.lanches:
            return None
        return max(self.lanches, key=lambda l: l.preco)
  \end{lstlisting}
\end{frame}

% ============================================================
% SLIDE 22
% ============================================================
\begin{frame}[fragile]{Usando o Cardápio}
  \begin{lstlisting}[language=Python]
cardapio = Cardapio("O Podrao")
cardapio.adicionar(Lanche("X-Tudo", 18.50, "pao, hamburguer..."))
cardapio.adicionar(Lanche("X-Salada", 14.00, "pao, alface..."))
cardapio.adicionar(Lanche("Misto Quente", 8.00, "pao, presunto"))

cardapio.mostrar()
  \end{lstlisting}

  Saída:
  \begin{lstlisting}[language=Python, numbers=none]
=== Cardapio: O Podrao ===
  X-Tudo - R$18.50
  X-Salada - R$14.00
  Misto Quente - R$8.00
  \end{lstlisting}
\end{frame}

% ============================================================
% SLIDE 23
% ============================================================
\begin{frame}[fragile]{Consultando o Cardápio}
  \begin{lstlisting}[language=Python]
print(f"Preco medio: R${cardapio.preco_medio():.2f}")
# Preco medio: R$13.50

campeao = cardapio.mais_caro()
print(f"Mais caro: {campeao.descrever()}")
# Mais caro: X-Tudo - R$18.50
  \end{lstlisting}

  \vspace{0.5cm}
  Note como \texttt{mais\_caro()} retorna um \textbf{objeto} \texttt{Lanche},
  e podemos chamar \texttt{descrever()} nele --- objetos conversam entre si!
\end{frame}

% ============================================================
% SLIDE 24
% ============================================================
\begin{frame}{Diagrama: Composição}
  % Diagram: Composição Cardapio-Lanche
  % Coordinates:
  %   cardapio at (0, 0) -- objeto Cardapio
  %   lanche1 at (6, 1.5) -- primeiro Lanche
  %   lanche2 at (6, 0) -- segundo Lanche
  %   lanche3 at (6, -1.5) -- terceiro Lanche
  \begin{center}
    \begin{tikzpicture}[
      classbox/.style={draw, rounded corners, minimum width=3.5cm,
                       minimum height=1.2cm, text width=3.3cm, align=center,
                       fill=blue!10},
      objbox/.style={draw, rounded corners, minimum width=3.0cm,
                     minimum height=0.9cm, text width=2.8cm, align=center,
                     fill=green!10},
      >=stealth
    ]
      \node[classbox] (cardapio) at (0, 0)
        {\textbf{Cardapio}\\``O Podrão''};
      \node[objbox] (l1) at (7, 1.5)
        {\footnotesize X-Tudo\\R\$18.50};
      \node[objbox] (l2) at (7, 0)
        {\footnotesize X-Salada\\R\$14.00};
      \node[objbox] (l3) at (7, -1.5)
        {\footnotesize Misto Quente\\R\$8.00};

      \draw[->, thick] (cardapio) -- (l1)
        node[midway, above, font=\footnotesize] {contém};
      \draw[->, thick] (cardapio) -- (l2);
      \draw[->, thick] (cardapio) -- (l3);

      \node[font=\footnotesize, gray] at (0, -2.2)
        {1 objeto Cardapio};
      \node[font=\footnotesize, gray] at (7, -2.2)
        {N objetos Lanche};
    \end{tikzpicture}
  \end{center}
\end{frame}

% ############################################################
% SEÇÃO 6: Resumo
% ############################################################
\section{Resumo}

% ============================================================
% SLIDE 25
% ============================================================
\begin{frame}{Resumo da Aula}
  \begin{table}
    \centering
    \small
    \begin{tabular}{lp{8cm}}
      \toprule
      \textbf{Conceito} & \textbf{O que aprendemos} \\
      \midrule
      Classe       & Modelo que define atributos e métodos \\
      Objeto       & Instância concreta com dados próprios \\
      \texttt{\_\_init\_\_} & Construtor --- inicializa cada objeto \\
      \texttt{self}  & Referência ao próprio objeto \\
      Atrib.\ de instância & Pertence ao objeto (criado em \texttt{\_\_init\_\_}) \\
      Atrib.\ de classe    & Compartilhado por todos os objetos \\
      Método       & Função que opera sobre o objeto \\
      Composição   & Objeto contém outros objetos (``tem-um'') \\
      \bottomrule
    \end{tabular}
  \end{table}
\end{frame}

% ============================================================
% SLIDE 26
% ============================================================
\begin{frame}{Armadilhas para Lembrar}
  \begin{enumerate}
    \item \textbf{Esquecer \texttt{self}} no primeiro parâmetro de métodos
          $\rightarrow$ \texttt{TypeError}
    \item \textbf{Listas/dicts como atributo de classe} $\rightarrow$
          compartilhamento indesejado entre objetos
    \item \textbf{Confundir classe com objeto} $\rightarrow$
          \texttt{Lanche.descrever()} sem instância dá erro
    \item \textbf{Não chamar \texttt{\_\_init\_\_} corretamente} $\rightarrow$
          atributos faltando (\texttt{AttributeError})
  \end{enumerate}
\end{frame}

% ============================================================
% SLIDE 27
% ============================================================
\begin{frame}{Próxima Aula: POO --- Parte 2}
  \begin{itemize}
    \item \textbf{Encapsulamento}: atributos públicos vs.\ privados
    \item \textbf{\texttt{@property}}: getters e setters pythônicos
    \item \textbf{Métodos especiais}: \texttt{\_\_str\_\_}, \texttt{\_\_repr\_\_},
          \texttt{\_\_eq\_\_}
    \item \textbf{Validação de dados}: garantir que preço $> 0$,
          nome não vazio
  \end{itemize}

  \vspace{0.5cm}
  O food truck ``O Podrão'' continua evoluindo!
\end{frame}

% ============================================================
% SLIDE 28 – Encerramento
% ============================================================
\begin{frame}[plain]
  \centering
  \vfill
  {\Large\textbf{Dúvidas?}}

  \vspace{1cm}
  {\normalsize Prof.\ Dr.\ André Yoshiaki Kashiwabara}\\
  {\small \texttt{andreyoshiaki@dcomp.ufs.br}}\\
  {\small DCOMP --- Universidade Federal de Sergipe}
  \vfill
\end{frame}

\end{document}
